% 原书第 157—159 页；导言及第 8.1 节全文。
\chapter{与机器学习的联系（二）}
\label{ch:machine-learning-two}

如第 5 章所述，统计模型背后的因果结构，可能对半监督学习、领域适应等机器学习任务产生重要影响。本章重新讨论这一主题，重点转向多元情形。首先介绍一种方法：对于给定的因果结构，利用机器学习对系统误差建模；接着讨论强化学习的一些问题，并介绍其在计算广告中的应用；最后讨论领域适应。

\section{半同胞回归}
\label{sec:ch08-half-sibling}

\textbf{半同胞回归（half-sibling regression）}利用给定的因果结构（见图 8.1），减少预测任务中的系统噪声。目标是重建未被观测到的信号 $Q$。Schölkopf et al.［2015］提出，可以从信号 $Y$ 中去除那些能够由其他测量值 $X$ 解释的信息，从而对 $Y$ 去噪；这里，$X$ 与 $Y$ 都受到同一噪声源的污染。具体而言，$X$ 是对某些信号 $R$ 的测量，而 $R$ 与 $Q$ 独立。直观地说，$Y$ 中能够由 $X$ 解释的所有成分，都必然来自系统噪声 $N$，因此应予去除。更准确地说，我们用
\[
 \widehat Q\coloneqq Y-\EE[Y\mid X]
\]
估计 $Q$。其中，$\EE[Y\mid X]$ 是将 $Y$ 对其“半同胞”$X$ 进行回归的结果；注意，$X$ 与 $Y$ 共有父节点 $N$，见图 8.1。

\begin{figure}[!htbp]
\centering
\begin{tikzpicture}[causal,x=1cm,y=1cm]
 \node[latent] (ch08f81Q) at (0,1.5) {$Q$};
 \node[latent] (ch08f81N) at (3.7,1.5) {$N$};
 \node[latent] (ch08f81R) at (7.4,1.5) {$R$};
 \node[vtx] (ch08f81Y) at (0,0) {$Y$};
 \node[vtx] (ch08f81X) at (7.4,0) {$X$};
 \draw[ed] (ch08f81Q)--(ch08f81Y);
 \draw[ed] (ch08f81N)--(ch08f81Y);
 \draw[ed] (ch08f81N)--(ch08f81X);
 \draw[ed] (ch08f81R)--(ch08f81X);
 \node[font=\small,anchor=east] at (-1.5,1.5) {未观测};
 \node[font=\small,anchor=east] at (-1.5,0) {已观测};
 \node[font=\small,above=4pt] at (ch08f81Q.north) {真实信号};
 \node[font=\small,above=4pt] at (ch08f81N.north) {系统噪声};
 \node[font=\small,above=4pt] at (ch08f81R.north) {其他信号};
 \node[font=\small,below=4pt] at (ch08f81Y.south) {关注的测量值};
 \node[font=\small,below=4pt] at (ch08f81X.south) {其他测量值};
\end{tikzpicture}
\caption{适用于搜寻系外行星问题的因果结构。底层的关注信号 $Q$ 只能以含噪版本 $Y$ 的形式测得。如果同一噪声源也污染了对其他信号的测量，而这些信号与 $Q$ 独立，就可以利用这些测量来去噪。本例中，望远镜 $N$ 构成系统噪声，影响对相互独立的光变曲线的测量 $X$ 和 $Y$。}
\label{fig:ch08-half-sibling}
\end{figure}


可以证明，对于任意满足 $Q\indep X$ 的随机变量 $Q,X,Y$，都有［Schölkopf et al., 2016，命题 1］
\[
 \EE\left[\bigl(Q-\EE[Q]-\widehat Q\bigr)^2\right]
 \leq \EE\left[\bigl(Q-\EE[Q]-(Y-\EE[Y])\bigr)^2\right].
\]
也就是说，这一方法的效果不会比直接采用测量值 $Y$ 更差。如果进一步假定系统噪声以加性的方式起作用，即存在某个未知函数 $f$，使得 $Y=Q+f(N)$，那么有［Schölkopf et al., 2016，命题 3］
\begin{equation}
 \EE\left[\bigl(Q-\EE[Q]-\widehat Q\bigr)^2\right]
 =\EE\bigl[\Var[f(N)\mid X]\bigr].
 \tag{8.1}
\end{equation}
如果加性噪声是 $X$ 的函数，即对于某个未知函数 $\psi$，有 $f(N)=\psi(X)$，那么式 (8.1) 的右侧为零。因此，$\widehat Q$ 能够恢复 $Q$，两者至多相差一个加性平移。其他充分条件见 Schölkopf et al.［2016］。

以搜寻系外行星为例。2009 年发射的 Kepler 太空望远镜，在银河系的一小片区域内搜寻系外行星，监测约 $150{,}000$ 颗恒星的亮度。\footnote{\url{https://en.wikipedia.org/wiki/Kepler_(spacecraft)}，访问日期：2016 年 7 月 13 日。}
如果某颗恒星周围存在行星，而且行星轨道合适，能够遮挡恒星的一部分，那么该恒星的光变曲线就会表现出周期性的光强下降，见图 8.2。这些测量受到望远镜所产生的系统噪声污染，使潜在的行星信号难以检测。幸运的是，望远镜会同时测量许多恒星。这些恒星彼此相隔数光年，可以假定它们在因果上相互独立，进而在统计上相互独立。因此，图 8.1 所示的因果结构非常适合这一问题，可以应用半同胞回归。这种简单方法的表现出人意料地好［Schölkopf et al., 2015］。

\begin{figure}[!htbp]
\centering
\begin{tikzpicture}[x=1cm,y=1cm]
 % Exact printed paths of original Figure 8.2; text labels translated separately.
\path[draw={rgb,1:red,0.01112383;green,0.01551843;blue,0.01744106},line width=0.46596pt] (12.844766,2.286337) ..controls (12.875814,2.718319) and (7.652225,2.758673)..(6.786787,2.746359) ..controls (5.921373,2.758673) and (0.718580,2.715446)..(0.728819,2.283476) ..controls (0.728819,2.725746) and (5.921373,2.758673)..(6.786787,2.758673) ..controls (7.652225,2.770999) and (12.886053,2.728607)..(12.844766,2.286337);
\path[draw={rgb,1:red,0.01112383;green,0.01551843;blue,0.01744106},line width=0.34799pt,fill={rgb,1:red,0.96623176;green,0.92373538;blue,0.07670710}] (8.017962,2.341030) ..controls (8.017962,3.011181) and (7.476532,3.554440)..(6.808676,3.554440) ..controls (6.140820,3.554440) and (5.599415,3.011181)..(5.599415,2.341030) ..controls (5.599415,1.670890) and (6.140820,1.127632)..(6.808676,1.127632) ..controls (7.476532,1.127632) and (8.017962,1.670890)..(8.017962,2.341030);
\path[draw={rgb,1:red,0.01112383;green,0.01551843;blue,0.01744106},line width=0.46596pt] (12.844791,2.283159) ..controls (12.875839,1.851189) and (7.652249,1.810823)..(6.786836,1.823149) ..controls (5.921398,1.810823) and (0.718605,1.854062)..(0.728856,2.286020) ..controls (0.728856,1.843762) and (5.921398,1.810823)..(6.786836,1.810823) ..controls (7.652249,1.798510) and (12.886077,1.840901)..(12.844791,2.283159);
\path[draw={rgb,1:red,0.01112383;green,0.01551843;blue,0.01744106},line width=0.34799pt,fill={rgb,1:red,0.93083084;green,0.12043946;blue,0.14123751}] (5.182437,1.823152) ..controls (5.182437,1.886045) and (5.131501,1.937031)..(5.068656,1.937031) ..controls (5.005824,1.937031) and (4.954888,1.886045)..(4.954888,1.823152) ..controls (4.954888,1.760246) and (5.005824,1.709260)..(5.068656,1.709260) ..controls (5.131501,1.709260) and (5.182437,1.760246)..(5.182437,1.823152);
\path[draw={rgb,1:red,0.01112383;green,0.01551843;blue,0.01744106},line width=0.34799pt,fill={rgb,1:red,0.93083084;green,0.12043946;blue,0.14123751}] (5.829365,1.822900) ..controls (5.829365,1.885806) and (5.778441,1.936791)..(5.715583,1.936791) ..controls (5.652739,1.936791) and (5.601815,1.885806)..(5.601815,1.822900) ..controls (5.601815,1.760006) and (5.652739,1.709021)..(5.715583,1.709021) ..controls (5.778441,1.709021) and (5.829365,1.760006)..(5.829365,1.822900);
\path[draw={rgb,1:red,0.01112383;green,0.01551843;blue,0.01744106},line width=0.34799pt,fill={rgb,1:red,0.93083084;green,0.12043946;blue,0.14123751}] (6.922452,1.823152) ..controls (6.922452,1.886045) and (6.871528,1.937031)..(6.808696,1.937031) ..controls (6.745851,1.937031) and (6.694927,1.886045)..(6.694927,1.823152) ..controls (6.694927,1.760246) and (6.745851,1.709260)..(6.808696,1.709260) ..controls (6.871528,1.709260) and (6.922452,1.760246)..(6.922452,1.823152);
\path[draw={rgb,1:red,0.01112383;green,0.01551843;blue,0.01744106},line width=0.59158pt] (0.708159,0.942702) --(5.611092,0.942702) ..controls (5.754239,0.942702) and (5.754239,0.594705)..(5.928225,0.572349) --(7.668272,0.572349) ..controls (7.842270,0.594705) and (7.842270,0.942702)..(8.024850,0.942702) --(12.888362,0.942702);
\path[draw={rgb,1:red,0.01112383;green,0.01551843;blue,0.01744106},line width=0.34799pt,fill={rgb,1:red,0.01112383;green,0.01551843;blue,0.01744106}] (0.399056,1.823141) --(0.509830,1.887176) --(0.509830,1.759105) --(0.399056,1.823141);
\path[draw={rgb,1:red,0.01112383;green,0.01551843;blue,0.01744106},line width=0.34799pt,fill={rgb,1:red,0.01112383;green,0.01551843;blue,0.01744106}] (12.972688,0.324357) --(12.861915,0.260309) --(12.861915,0.388392) --(12.972688,0.324357);
\path[draw={rgb,1:red,0.01112383;green,0.01551843;blue,0.01744106},line width=0.69598pt,fill={rgb,1:red,0.01112383;green,0.01551843;blue,0.01744106}] (0.454446,0.323470) --(0.454446,1.822916);
\path[draw={rgb,1:red,0.01112383;green,0.01551843;blue,0.01744106},line width=0.69598pt,fill={rgb,1:red,0.01112383;green,0.01551843;blue,0.01744106}] (12.845007,0.323470) --(0.454441,0.323470);

 \node[font=\small,anchor=west] at (7.96,3.2) {恒星};
 \node[font=\small] at (4.2,2.13) {行星};
 \node[font=\small,rotate=90] at (.17,1.3) {亮度};
 \node[font=\small] at (10.4,.65) {光变曲线};
 \node[font=\small,anchor=north east] at (12.8,.16) {时间};
\end{tikzpicture}
\caption{每当行星遮挡恒星的一部分，光强就会下降。如果行星绕恒星运行，这一现象便会周期性发生。（图片由 Nikola Smolenski 提供，\url{https://en.wikipedia.org/wiki/File:Planetary_transit.svg}，许可协议 CC BY-SA 3.0。为使图示更清晰并统一风格，原书对图片作了编辑。）}
\label{fig:ch08-transit}
\end{figure}


其他应用领域也曾使用相关方法，但并未诉诸因果建模［Gagnon-Bartsch and Speed, 2012；Jacob et al., 2016］。一旦考虑问题的因果结构（图 8.1），上述方法便自然浮现出来，而且还能据此给出支持该方法的理论论证。
\FloatBarrier
